\documentclass[aspectratio=169]{beamer}
\usepackage[utf8]{inputenc}
\usepackage[russian]{babel}
\usepackage{listings}
\usepackage{xcolor}
\usepackage{hyperref}

\usetheme{Madrid}
\usecolortheme{default}

\definecolor{codegray}{rgb}{0.5,0.5,0.5}
\definecolor{backcolour}{rgb}{0.95,0.95,0.92}

\lstdefinestyle{mystyle}{
    backgroundcolor=\color{backcolour},   
    basicstyle=\ttfamily\scriptsize,
    breaklines=true,
    xleftmargin=2em,
    framexleftmargin=1.5em,
}

\lstset{style=mystyle}

\title{Claude Skills}
\subtitle{Формат записи повторяющихся инструкций для LLM}
\date{17 февраля 2026}

\begin{document}

\frame{\titlepage}

% ==================== СЛАЙД 1: ПРОБЛЕМА ====================

\begin{frame}
\frametitle{Проблема: объяснять одно и то же}

\textbf{LLM не всё знает и умеет — нужны инструкции двух типов:}

\vspace{0.3cm}

\textbf{Требования (ЧТО делать):}
\begin{itemize}
    \item «Используй Stripe API v2024-10, не v2023-08. Всегда webhook-подтверждение»
    \item «Отчёты клиентам: сначала executive summary, потом метрики, потом рекомендации»
\end{itemize}

\vspace{0.3cm}

\textbf{Обучающие инструкции (КАК делать):}
\begin{itemize}
    \item «Для Excel с формулами: используй openpyxl, пиши формулы через .value = '=SUM(A1:A10)', не через numbers»
    \item «Statistical significance в A/B тестах: chi-square для категориальных данных, t-test для числовых. Формула chi-square: ...»
    \item «PostgreSQL async/await: всегда используй async with pool.acquire(), закрывай connection в finally блоке»
\end{itemize}

\vspace{0.3cm}

\begin{block}{Claude Skills = формат записи обоих типов инструкций}
\end{block}

\end{frame}

% ==================== СЛАЙД 2: РЕШЕНИЕ ====================

\begin{frame}
\frametitle{Решение: Claude Skills}

\textbf{Skills} = папка с инструкциями, которую Claude подгружает автоматически

\vspace{0.3cm}

\begin{block}{Как работает}
\begin{enumerate}
    \item Один раз описываешь процесс в файле \texttt{SKILL.md}
    \item Claude сам определяет, когда применить skill
    \item Динамически подгружает только нужные инструкции
    \item Выполняет задачу по стандарту без повторных объяснений
\end{enumerate}
\end{block}

\vspace{0.3cm}

\textbf{Результат:} Консистентность + экономия времени + масштабируемость

\end{frame}

% ==================== СЛАЙД 3: СТРУКТУРА ====================

\begin{frame}
\frametitle{Структура skill: это просто папка}

\textbf{Базовая структура:}

\vspace{0.3cm}

\begin{itemize}
    \item \texttt{telegram-post-writer/}
    \begin{itemize}
        \item \texttt{SKILL.md} — главный файл с правилами
        \item \texttt{instructions.md} — детальные инструкции (опционально)
        \item \texttt{examples/} — примеры входов и выходов (опционально)
        \item \texttt{scripts/} — Python/bash для автоматизации (опционально)
    \end{itemize}
\end{itemize}

\vspace{0.5cm}

\begin{block}{SKILL.md — минимальная структура}
\begin{itemize}
    \item \texttt{name:} название skill
    \item \texttt{description:} что делает
    \item Цель, правила, стоп-слова, формат выхода
\end{itemize}
\end{block}

\end{frame}

% ==================== СЛАЙД 4: ПРИМЕР ====================

\begin{frame}
\frametitle{Лучший пример для нетехнических людей: Email Response}

\begin{block}{Задача}
Отвечать на письма клиентов в едином стиле: вежливо, коротко, по делу
\end{block}

\vspace{0.3cm}

\textbf{Было (без skill):}

Каждый раз объясняешь: «Будь вежлив, но не формален. Используй имя клиента. Сначала извинись, потом решение, потом предложи помощь»

\vspace{0.5cm}

\textbf{Стало (с skill):}

Создал skill «customer-support-email» один раз. Теперь просто пишешь: «Ответь на это письмо», и Claude сам применяет все правила:

\begin{itemize}
    \item Обращается по имени
    \item Признаёт проблему клиента
    \item Даёт конкретное решение за 2-3 предложения
    \item Предлагает дополнительную помощь
    \item Подписывается твоим именем
\end{itemize}

\end{frame}

% ==================== СЛАЙД 5: АВТОМАТИЗМ ====================

\begin{frame}
\frametitle{Почему это не просто папка с файлами}

\textbf{Два ключевых отличия:}

\vspace{0.5cm}

\textbf{1. Готовые skills можно скачать}

Есть хабы (agentskills.io, agentskillshub.dev), где тысячи готовых skills под разные задачи. Не нужно писать с нуля — берёшь готовый и адаптируешь.

\vspace{0.5cm}

\textbf{2. Агенты сами понимают, когда использовать skill}

Это уже заложено в многих агентах (Claude Code, Claude Desktop). Ты просто пишешь запрос, агент анализирует контекст и автоматически подгружает нужный skill без явной команды.

\vspace{0.5cm}

\begin{alertblock}{Главное}
Skill = стандартизированный формат, который понимают разные агенты
\end{alertblock}

\end{frame}

% ==================== СЛАЙД 6: ДЕПЛОЙ ПРИМЕР ====================

\begin{frame}
\frametitle{Практический пример: деплой на сервер}

\textbf{Структура deploy-skill:}

\vspace{0.3cm}

\begin{itemize}
    \item \texttt{deploy-skill/}
    \begin{itemize}
        \item \texttt{SKILL.md} — описание процесса деплоя
        \item \texttt{scripts/deploy.sh} — rsync + restart service
        \item \texttt{scripts/backup.sh} — автобэкап перед деплоем
        \item \texttt{scripts/rollback.sh} — откат при ошибке
        \item \texttt{examples/nginx.conf} — шаблон конфига
    \end{itemize}
\end{itemize}

\vspace{0.5cm}

\textbf{Запрос:} «Задеплой проект на прод»

\textbf{Claude автоматически:}
\begin{itemize}
    \item Создаёт бэкап текущей версии
    \item Загружает новую версию через rsync
    \item Перезапускает сервис и проверяет статус
\end{itemize}

\end{frame}

% ==================== СЛАЙД 7: ХАБЫ ====================

\begin{frame}
\frametitle{Где найти готовые skills}

\begin{block}{Основные хабы}
\begin{itemize}
    \item \textbf{agentskills.io} — официальный стандарт и каталог
    \item \textbf{agentskillshub.dev} — проверенные skills с security-сканированием
    \item \textbf{claude.ai} — встроенный каталог (Pro/Max/Team/Enterprise)
    \item \textbf{GitHub} — community-репозитории
\end{itemize}
\end{block}

\vspace{0.3cm}

\textbf{Примеры skills из хабов:}
\begin{itemize}
    \item \textbf{Notion integration} — экспорт задач и заметок
    \item \textbf{Code review assistant} — анализ PR по стандартам команды
    \item \textbf{SEO content generator} — статьи с ключевыми словами
    \item \textbf{Expense audit} — проверка счетов и синхронизация с QuickBooks
\end{itemize}

\end{frame}

% ==================== СЛАЙД 8: ОПАСНОСТЬ ====================

\begin{frame}
\frametitle{Опасность: Prompt Injection}

\begin{alertblock}{Реальная атака: «What Would Elon Do?»}
Skill из топ-1 OpenClaw репозитория оказался \textbf{вредоносным}
\end{alertblock}

\vspace{0.3cm}

\textbf{Что делал skill:}
\begin{itemize}
    \item Exfiltration данных на сервер атакующего через \texttt{curl}
    \item Prompt injection для обхода safety guidelines
    \item Выполнение команд без подтверждения пользователя
\end{itemize}

\vspace{0.3cm}

\begin{alertblock}{Ключевая проблема}
Skill исполняется с \textbf{теми же правами, что и сам агент}. \\
Команды агента \textbf{по-сути никто не читает} перед разрешением.
\end{alertblock}

\vspace{0.3cm}

\textbf{Источники:} 
\small{\url{https://www.authmind.com/post/openclaw-malicious-skills-agentic-ai-supply-chain}}


\end{frame}

% ==================== СЛАЙД 9: ВЫВОДЫ ====================

\begin{frame}
\frametitle{Ключевые выводы}

\begin{enumerate}
    \item \textbf{Проблема решена} — Skills избавляют от повторных объяснений
    
    \item \textbf{Автоматизм} — Claude сам выбирает и применяет skills
    
    \item \textbf{Практичность} — от контента до деплоя, реальные кейсы
    
    \item \textbf{Опасность} — prompt injection и malware в community-skills
    
    \item \textbf{Критичная проблема безопасности} — skill выполняется с правами агента, команды не проверяются пользователем
\end{enumerate}

\vspace{0.5cm}

\begin{block}{Источники}
\begin{itemize}
    \item Habr: \url{https://habr.com/ru/articles/993174/}
    \item Роадмап: \url{https://t.me/NeuralProfit/2143}
    \item Хабы: \url{https://agentskills.io}, \url{https://agentskillshub.dev}
\end{itemize}
\end{block}

\end{frame}

\end{document}
